% Draft-II introduction: paper-level architecture centered on the planar theory.

\section{Introduction}\label{sec:intro}

\subsection{The routing problem}

Let $s\ge1$, and let $A_1,\dots,A_s$ be pairwise disjoint full continua in a Jordan disk
$U\subset\C$.  The elementary task of joining the $A_i$ becomes substantially
less elementary when the connector must also avoid a frozen compact obstacle
set and preserve designated prime ends of the complements
$\Chat\setminus A_i$.  A connector can erase such an access by entering its
selected side, identify two accesses that were distinct before enlargement,
or disconnect the region through which later ports were meant to run.
None of these failures is detected by connectedness of the final union.

This problem appears naturally when small filled Julia sets, puzzle pieces,
or combinatorial fibers must be organized without assuming local
connectivity.  Polynomial-like maps provide the basic continua
\cite{DH85}; Kahn--Lyubich impose finite-decoration combinatorics and
prove a~priori bounds for that class \cite{KL08}; recent work on
primitive renormalization constructs native or model-lifted Hubbard trees
when the required dynamical structure is available \cite{KKL26}; and
Dudko's roadmap emphasizes connected hulls of families of small Julia sets
as one ingredient in the broader local-connectivity program \cite{Dud25}.
The present paper isolates the planar problem that remains after the
dynamical data have selected finitely many continua and accesses but before
any canonical connector has been justified.

The hypotheses deliberately allow non-locally-connected boundaries.
A protected access is therefore not represented by a presumed landing ray
or boundary parametrization.  It is represented by a null prime-end
crosscut chain together with its selected sides.  The part forbidden to a
future connector is the \emph{relative collar tail}
$\overline D\setminus A_i$, not the whole closed side $\overline D$.
This distinction is essential: different prime ends may have overlapping
impressions inside the same atom even though their relative collar tails in
the complement are disjoint.

\subsection{Main results}

The first conclusion is negative but foundational.  There is no least
compact connected full hull even for two points
(Theorem~\ref{thm:no-least}), so ``take the smallest carrier'' is not a
definition.  Nor can one safely choose a carrier and absorb pulled-back
excess monotonically: Proposition~\ref{prop:access-collision} gives a
degree-one conformal branch, with no critical-value obstruction, in which
the first absorption erases a protected access permanently.  These examples
show that the isotopy/access class of the routing is additional data.

The positive replacement is a \emph{protected ribbon-routing certificate}.
It records free singleton ports from the atoms, a finite embedded tree
joining their outer endpoints, the induced ribbon order, every protected
crosscut chain, and the complete frozen inventory to be avoided.
Theorem~\ref{thm:realization} realizes any such certificate by a compact
connected full carrier.  Every protected chain remains a prime-end chain of
the enlarged complement, distinct protected chains remain distinct whenever
their relative collar tails are disjoint, and cyclic order is preserved
within each common source complement.  No local connectivity of an atom is
used.

The first structural result developed here is an exact attachment ledger.
For a compact connected connector $C$, full atoms $A_i$, and traces
$T_i=C\cap A_i$, Theorem~\ref{thm:separation-budget} extracts from the
compact-pair Mayer--Vietoris sequence a natural short exact sequence at
every attachment.  When the relevant component counts are finite, it gives
\[
 \sigma\!\left(C\cup\bigcup_iA_i\right)
 =
 \sigma(C)+
 \sum_i\bigl(\#\pi_0(T_i)-1-\sigma(T_i)\bigr),
\]
where $\sigma(K)=\#\pi_0(\Chat\setminus K)-1$.  The formula is a
credit/debit law.  Each extra component of an attachment trace creates a
complementary component; each complementary component of the trace fills
one separation already present in the connector.  In particular, if $C$
and the atoms are full, every $T_i$ is nonseparating and the union is full
if and only if all traces are connected
(Corollary~\ref{cor:full-connector-traces}).  Thus the familiar tree with
singleton attachments is only one sufficient model: an arbitrary full
continuum connector with arbitrary connected traces works just as well.

The same ledger identifies a compactness gate for this framework.
Proposition~\ref{prop:trace-birth} constructs, for every $m$, full arc
connectors $C_n$ meeting an atom once and converging in the Hausdorff metric
to a full arc $C$ whose trace has $m+1$ components.  Each finite-stage
union $A\cup C_n$ is full, while the
limit union has exactly $m$ bounded complementary components.
Connector cycles can likewise be born in a limit.  Consequently neither
fullness nor the singleton-trace condition passes to a virtual or infinite
carrier by Hausdorff compactness alone: a limiting argument must separately
control connector cohomology and trace multiplication
(Warning~\ref{warn:compactness-gate}).  Theorem
\ref{thm:nested-limit-stability} gives a positive counterpart: a decreasing
sequence of full connector envelopes with connected nonempty traces has a
full limiting connector with connected traces, and common protected data
survive.  This identifies a concrete inverse-limit architecture that avoids
both defect-birth mechanisms.

The second structural result is an exact routing criterion.  Freeze the
atoms, relative protected collar tails, obstacles, critical exclusions, and
localization-forbidden data into a compact set $Q$, and put
$\Omfree=U\setminus Q$.  Theorem~\ref{thm:free-component} says that a
protected certificate exists if and only if one component of $\Omfree$
admits simultaneous singleton ports to every required atom.  If no
component of this fixed free domain has simultaneous ports, the failure
persists under every connector enlargement respecting the same $Q$.  The
proof uses a one-ended slit lemma whose statement keeps
the free tips inside the residual domain.  The source version deleted those
tips and mapped a boundary-anchored slit into an open domain; the corrected
Lemma~\ref{lem:slits} repairs both defects, and the proof of the criterion is
rebuilt around the repaired statement.

Finally, the exact criterion is converted into finite audits.
Theorem~\ref{thm:common-domain} packages the frozen inventory into pairwise
disjoint full continua and uses compacta Alexander duality to force a common
free domain.  Theorem~\ref{thm:audit} proves that fullification is independent
of processing order for a fixed typed assignment and reduces success to
coverage, packet separation/localization, and a nonempty unprotected
exposure locus at each atom.  Packet merger, localization escape, and port
loss come with witnesses persistent under every same-label monotone
enlargement.  Proposition~\ref{prop:cellwise} supplies a directly checkable
Jordan-cell criterion.

\subsection{Position}

The underlying tools are classical: Janiszewski--\allowbreak Mullikin separation
\cite{Jan15,Mul22,New51,Why42}, the classical theories of unicoherence and
multicoherence \cite{Bor31,Eil36}, Carath\'eodory prime ends \cite{Pom92}, the
Jordan--Schoenflies theorem, filled hulls, and compacta Alexander duality and
Mayer--Vietoris \cite[Chapter~6]{Spa66}.  We do not claim a new foundation
for planar topology.  The contribution is a theorem-level synthesis aimed
at a specific proof obligation that the standard constructions do not
record: finite connection with named access survival, exact existence and
failure tests, and a quantitative accounting of the separations created or
destroyed by attachment.

In particular, the separation-ledger theorem is formally a consequence of
classical duality and Mayer--Vietoris; its full/full zero-defect criterion
is part of the Janiszewski--Mullikin tradition.  Its value here is the
homomorphism-level diagnostic, the creation-versus-cancellation formula,
and the limit obstruction it exposes.  The other central statements
turn prime-end intentions into finite data: null chains and relative
collar tails, singleton ports, an embedded ribbon tree, a frozen free domain,
typed packet assignments, and persistent witnesses.  Section
\ref{sec:comparison} compares this package with least hulls, monotone
absorption, regulated and Hubbard trees, Moore quotients, and a
lamination-lift candidate.

This architecture is intentionally distinct from the companion
direct-decoration manuscript \cite{Wes26}.  That manuscript contains the rooted-end
compactification, partial-shift dynamics, gap cells, arbitrary packet
replacement, and the finite decoration-carrier application.  The present
paper contains the general planar certification theory, the exact
free-component and packet criteria, the attachment ledger, and the
Hausdorff-limit no-go theorem.  Separating the two removes a long duplicated
application and makes the logical dependency transparent.

\subsection{Scope and organization}

Everything here is finite and planar.  The paper does not construct a
virtual carrier or ``Wave horizon'' (the source manuscript's proposed
infinite limiting carrier), an analytic correspondence, a uniform
incidence bound, or a compression mechanism, and it proves no new case of
local connectivity of the Mandelbrot set.  The recent theorem of
Kahn--Kapiamba--Lyubich applies in its stated parabolically bounded primitive
category \cite{KKL26}; it is a positive benchmark, not a consequence of the
present framework.  Our conclusion is conditional in exactly one transparent
sense: once finite protected data or a valid packet assignment are supplied,
the planar realization problem is solved, while the production of coherent
data at an infinite dynamical horizon remains open.

Section~\ref{sec:prelim} fixes the compacta and prime-end conventions.
Section~\ref{sec:obstructions} proves the two motivating obstructions, and
Section~\ref{sec:comparison} places the framework among standard
constructions.  Section~\ref{sec:certificates} gives certificate realization
and prime-end survival.  Section~\ref{sec:separation-budget} proves the
attachment ledger and its sharp Hausdorff-limit obstruction.
Section~\ref{sec:criterion} proves the corrected slit lemma and exact
free-component criterion.  Sections~\ref{sec:packets} and~\ref{sec:audit}
develop the packet criterion, order-independent fullification, persistent
audits, and the cellwise test.  Section~\ref{sec:scope} closes with the exact
interface to future finite and limiting applications.
